\newpage
\begin{center}
    \phantomsection
    \addcontentsline{toc}{section}{RINGKASAN}
    \section*{\textbf{RINGKASAN}}
\end{center}

Korupsi masih menjadi permasalahan serius di Indonesia dan tidak hanya berkaitan dengan lemahnya penegakan hukum, tetapi juga berhubungan dengan nilai moral, religiusitas, dan perilaku sosial masyarakat. Indonesia dikenal sebagai negara dengan tingkat religiusitas yang tinggi, namun praktik korupsi masih terus terjadi di berbagai sektor. Berdasarkan hal tersebut, kegiatan ini dilakukan untuk memahami pandangan agama dalam menanamkan moral antikorupsi kepada masyarakat, khususnya dalam melihat hubungan antara ajaran agama, nilai integritas, dan perilaku antikorupsi dalam kehidupan sehari-hari. Kegiatan ini menggunakan pendekatan \textit{mixed-methods}, yaitu pengumpulan data kuantitatif melalui kuesioner kepada 177 responden dan data kualitatif melalui wawancara dengan lima tokoh agama dari Buddha, Kristen, Katolik, Hindu, dan Islam. Hasil kegiatan menunjukkan bahwa 95,5\% responden mengetahui bahwa agama mereka secara eksplisit melarang korupsi dan 94,4\% responden mengetahui adanya sanksi moral atau spiritual bagi pelaku korupsi. Namun, dampak ceramah antikorupsi terhadap perilaku responden hanya memperoleh skor rata-rata 3,88 dari skala 5, sedangkan pengetahuan responden terhadap tokoh agama yang aktif mengampanyekan antikorupsi memperoleh skor rata-rata 3,67. Hasil wawancara menunjukkan bahwa seluruh agama memandang korupsi sebagai pelanggaran moral yang serius karena berkaitan dengan tindakan mengambil hak orang lain secara tidak sah. Oleh karena itu, penanaman moral antikorupsi perlu diperkuat melalui pendidikan agama yang lebih kontekstual, keteladanan tokoh agama, serta pembiasaan nilai kejujuran dan integritas dalam masyarakat.

\noindent\textbf{Kata kunci :} nilai keagamaan; antikorupsi; moral; multireligius; integritas